\subsection{Greedy Algorithm Design}

The algorithm uses the earliest-finish-time greedy
strategy. It sorts the candidate windows in ascending
order of finishing time, then scans them in that order,
selecting a window if it starts no earlier than the
finish of the previously selected window and skipping
it otherwise.

\subsubsection{Algorithm Pseudocode}

\begin{verbatim}
GREEDY-WINDOW-SCHEDULING(W)
Input:  W = candidate green-light windows
Output: S = maximum set of compatible windows

1. Sort W by ascending finishing time
2. S = empty list;  lastFinish = -infinity
3. for each window w in W:
4.     if w.start >= lastFinish:
5.         add w to S;  lastFinish = w.finish
6. return S
\end{verbatim}

\subsubsection{Greedy Choice}

The greedy choice is the compatible window with the
earliest finishing time: finishing earlier leaves at
least as much time for later windows as any other
choice. Its optimality is proved by an exchange
argument in Section~\ref{sec:greedy-proof}.

\subsubsection{Explanation in the Traffic Signal Domain}

Each candidate window is a period in which a traffic
movement could receive a green signal. The scheduling
system arranges these opportunities by ending time,
approves the earliest-ending one, and then approves
each subsequent opportunity only if it does not overlap
an already approved green-light window.

For example, consider the following windows:

\begin{center}
\begin{tabular}{|c|c|c|}
\hline
Window & Start (s) & Finish (s) \\
\hline
1 & 0 & 10 \\
2 & 2 & 5 \\
3 & 5 & 8 \\
4 & 7 & 12 \\
5 & 8 & 14 \\
\hline
\end{tabular}
\end{center}

After sorting by finishing time, the algorithm selects
window 2 (finishing at 5\,s), window 3 (starting at
5\,s), and window 5 (starting at 8\,s, when window 3
finishes): three compatible green-light windows.
This optimizes the number of selected windows under the
simplified model; it does not directly optimize vehicle
throughput, signal safety, or network-wide congestion.
